% Pursuing Stacks (cotes 134-2 to 134-8): archive cote 134-7, site page 10 (PDF page 11), typescript page 452.
% typescript. Draft compared: ps6.tex:430-509.
% AI-transcribed by Claude Opus 5.5 (claude-opus-5-5, adaptive thinking, effort high), Message Batches API, batch ps3 (msgbatch_01XdFmuVVBAcgc7EFxdtCqsj), 2026-11-12; not yet checked by eye.
\scanpage{134-7/10}
\mspage{452}

\hand{(}p') \qquad $\Gamma(f) \ : \ M \to M'$

(corresponding to the action of f, viewed as a map $W(M)\to W(M')$, on
\uline{sections} of $W(M)$) - and likewise after any ring extension $k\to k'$ ,
defining a map
\[ \Gamma(f)_{k'} \ : \ M\otimes_k k' \to M'\otimes_k k' \]
from $W(M)$ to $W(M')$ - which is the description of the map of sheaves
$W(M)\longrightarrow W(M')$ associated to a map \hand{(}p). We have thus a description, in
terms of linear algebra, of a category $U(k)$, and of a \inserted{``sections''} functor

\xout{\unsure{hx}}(r) \qquad $\hand{\Gamma} \ : \ U(k) \longrightarrow (Sets) \quad , \qquad X \mapsto X(k)$ ,

which is essentially the functor \hand{(}h) above, viewed in a different light.
Now to our

\uline{Question}: \margin{\unsure{$\#$}}The category $U(k)$ and the functor r) being defined as above
in terms of linear algebra over the ground ring k , let $X_{\hand{*}} = (n\mapsto X_n)$
be a semisimplicial object in $U(k)$, consider the corresponding semi-sim-%
plicial set $X_{\hand{*}}(k) = (n\mapsto X_n(k))$, and the \unsure{semi-simplicial|semi-sipplicial} k-module
\margin{* \unsure{One has to} take $k=\underline{Z}$, and components \unsure{$X_n$|to be} projective.}
$L(X_{\hand{*}}) = (n\mapsto L(X_n))$ , where for an object M of $U(k)$, $L(M)$ is defined as
\[ \text{(s)} \qquad L(M) = \Gamma_k{}^{\wedge}(M) \quad , \]
which depends functorially on M in $U(k)$. Then in the derived category
of (Ab), is there a canonical isomorphism between $L(X_{\hand{*}})$ (with ground
ring restricted from k to $\underline{Z}$ ) and the abelianization $\underline{Z}^{(X_{\hand{*}}(k))}$ of $X_{\hand{*}}(k)$?

We may have to throw in some extra assumption on $X_{\hand{*}}$ , at any rate
\[ \hand{(}\text{t)} \qquad X_o = e \deleted{,} \hand{(final object of $U(k)$),} \]
\deleted{(the final object of U(k)),} giving rise to \xout{the} $L(X_o)\hand{\simeq} k$ . Also we may
have to restrict to $k=\underline{Z}$ , or otherwise correct the obvious drawback that
the two chain complexes do'nt have isomorphic $H_o$ (one is k I guess,
a k-module in any case, the other is $\underline{Z}$ ), by truncating accordingly
\margin{X}the two chain complexes ("killing" their $H_o$). The\over{\unsure{re}}{re} is a natural can-%
didate \over{\unsure{for}}{for} a map
\[ \hand{(}\text{u)} \qquad \underline{Z}^{(X_{\hand{*}}(k))} \longrightarrow L(X_{\hand{*}}) \quad , \]
by using the functorial map
\[ \hand{(}\text{u')} \qquad \underline{Z}^{(M)} \longrightarrow L(M) = \Gamma_k{}^{\wedge}(M) \quad , \]
deduced from the inclusion \hand{(n)}
\[ M \hookrightarrow L(M) \quad , \]
and we may still specify the question above, by asking \uline{whether} \hand{(}u) \uline{induces}
\uline{an isomorphism for homology groups} in dimension $i>0$.

\margin{10}
